\documentclass[11pt,a4paper]{article}

% --- PACKAGES ---
\usepackage[utf8]{inputenc}
\usepackage[indonesian]{babel}
\usepackage{amsmath, amssymb, amsfonts}
\usepackage{geometry}
\geometry{margin=1in, headheight=14pt}
\usepackage{graphicx}
\usepackage{booktabs}
\usepackage{tcolorbox}
\tcbuselibrary{skins,breakable}
\usepackage{enumitem}
\usepackage{fancyhdr}

% --- STYLING & COLORS ---
\definecolor{unibblue}{RGB}{0, 51, 102}
\definecolor{unibgold}{RGB}{212, 175, 55}
\definecolor{darkgreen}{RGB}{34, 139, 34}

\pagestyle{fancy}
\fancyhf{}
\rhead{\textbf{TKE-41XX: Optimisasi untuk Teknik Elektro}}
\lhead{Lembar Kerja Mahasiswa (LKM) Minggu 14}
\rfoot{Halaman \thepage}

\title{\vspace{-1cm}\textbf{LEMBAR KERJA MAHASISWA (LKM) MINGGU KE-14:}\\Optimisasi Multi-Objektif dan Metode Hibrid (Sintesis EE dan Mesin)}
\author{Program Studi S1 Teknik Elektro --- Universitas Bengkulu}
\date{Semester Ganjil 2026}

\begin{document}

\maketitle

\begin{tcolorbox}[colback=unibblue!5,colframe=unibblue,title=\textbf{Identitas Mahasiswa dan Instruksi}]
\begin{tabular}{ll}
\textbf{Nama Mahasiswa} : & \hrulefill\hspace{6cm} \\
\textbf{NPM} : & \hrulefill\hspace{6cm} \\
\textbf{Kelompok / Kelas} : & \hrulefill\hspace{6cm} \\
\end{tabular}
\vspace{0.2cm}
\begin{itemize}[leftmargin=*]
    \item Kerjakan LKM ini secara kolaboratif dalam kelompok (2--3 mahasiswa).
    \item Jawablah setiap pertanyaan secara sistematis pada ruang yang disediakan.
    \item Soal ini mencakup ranah kognitif Taksonomi Bloom tingkat **C1 hingga C6**.
\end{itemize}
\end{tcolorbox}

\vspace{0.3cm}

\section*{Bagian 1: Konsep Dasar Multi-Objektif dan Dominasi Pareto (C1--C2)}

\begin{enumerate}[label=\textbf{Soal \arabic*:}, leftmargin=*]
    \item \textbf{[C1 --- Mengingat]} Jelaskan perbedaan utama antara optimisasi *Single-Objective* dan *Multi-Objective* (MOO) dalam konteks pengoperasian sistem tenaga listrik modern (*Combined Economic dan Emission Dispatch*)!
    
    \begin{tcolorbox}[colback=white,colframe=gray!50,height=3cm,title=Ruang Jawaban Soal 1]
    \end{tcolorbox}

    \item \textbf{[C2 --- Memahami]} Berikan kriteria formal matematis agar suatu alokasi daya $\mathbf{P}_a$ dikatakan **mendomina secara Pareto** terhadap alokasi daya $\mathbf{P}_b$ ($\mathbf{P}_a \succ \mathbf{P}_b$) untuk dua fungsi tujuan (Biaya dan Emisi)!
    
    \begin{tcolorbox}[colback=white,colframe=gray!50,height=3.5cm,title=Ruang Jawaban Soal 2]
    \end{tcolorbox}
\end{enumerate}

\section*{Bagian 2: Aplikasi Deterministik dan Sintesis Teknik Mesin (C3--C4)}

\begin{enumerate}[label=\textbf{Soal \arabic*:}, leftmargin=*, resume]
    \item \textbf{[C3 --- Mengaplikasikan]} Tiga buah solusi kandidat alokasi daya pada sistem generator menghasilkan nilai biaya $F_{\text{cost}}$ (\$/jam) dan emisi $E_{\text{emission}}$ (kg/jam) sebagai berikut:
    \begin{itemize}
        \item Solusi A: $(8200, 510)$
        \item Solusi B: $(8350, 440)$
        \item Solusi C: $(8280, 480)$
        \item Solusi D: $(8400, 520)$
    \end{itemize}
    Tentukan manakah di antara keempat solusi tersebut yang merupakan **solusi non-terdominasi** (*Pareto Optimal Set*) dan gambarkan batas *Pareto Front* secara konseptual!
    
    \begin{tcolorbox}[colback=white,colframe=gray!50,height=5cm,title=Ruang Jawaban Soal 3]
    \end{tcolorbox}

    \item \textbf{[C4 --- Menganalisis]} Bandingkan peran algoritma *Quasi-Newton BFGS* (dari materi Teknik Mesin Bab 6) dengan algoritma *Augmented Lagrangian* (dari materi Teknik Mesin Bab 7) jika diintegrasikan dengan PSO pada penyelesaian masalah CEED! Mengapa pendekatan hibrid ini lebih unggul dibandingkan PSO murni?
    
    \begin{tcolorbox}[colback=white,colframe=gray!50,height=5.5cm,title=Ruang Jawaban Soal 4]
    \end{tcolorbox}
\end{enumerate}

\section*{Bagian 3: Evaluasi dan Perancangan Algoritma Hibrid (C5--C6)}

\begin{enumerate}[label=\textbf{Soal \arabic*:}, leftmargin=*, resume]
    \item \textbf{[C5 --- Mengevaluasi]} Suatu *Pareto Front* CEED memiliki rentang biaya $F_{\text{cost}} \in [8194, 8342]$ dan emisi $E_{\text{emission}} \in [442, 509]$. Jika Solusi Kompromi C memiliki nilai $(8245, 465)$, hitunglah nilai derajat keanggotaan *Fuzzy* $\mu_{\text{cost}}$ dan $\mu_{\text{emission}}$ serta nilai kelayakan total $\mu^{(C)}$! Evaluasi apakah solusi ini layak dipilih sebagai *Best Compromise Solution*.
    
    \begin{tcolorbox}[colback=white,colframe=gray!50,height=4.5cm,title=Ruang Jawaban Soal 5]
    \end{tcolorbox}

    \item \textbf{[C6 --- Mencipta]} Rancanglah skema arsitektur algoritma optimisasi hibrid **BFGS-MOPSO** yang menggabungkan kemampuan eksplorasi global partikel MOPSO dengan kemampuan akselerasi lokal *Quasi-Newton BFGS*. Buatlah diagram alir (*flowchart*) atau *pseudocode* lengkapnya!
    
    \begin{tcolorbox}[colback=white,colframe=gray!50,height=6.5cm,title=Ruang Jawaban Soal 6]
    \end{tcolorbox}
\end{enumerate}

\end{document}
